% Generated draft: requires mathematical and source review.
\documentclass[11pt]{article}
\usepackage[margin=25mm]{geometry}
\usepackage{fontspec}
\setmainfont{Latin Modern Roman}
\usepackage{amsmath,amssymb,mathtools,unicode-math}
\setmathfont{Latin Modern Math}
\usepackage{enumitem,needspace,xurl,hyperref}
\setlength{\parindent}{0pt}
\setlength{\parskip}{5pt}
\setlength{\emergencystretch}{3em}
\newcommand{\N}{\mathbb{N}}
\newcommand{\Z}{\mathbb{Z}}
\newcommand{\Q}{\mathbb{Q}}
\newcommand{\R}{\mathbb{R}}
\newcommand{\C}{\mathbb{C}}
\newcommand{\F}{\mathbb{F}}
\newcommand{\A}{\mathbb{A}}
\newcommand{\PP}{\mathbb{P}}
\newcommand{\E}{\mathbb{E}}
\newcommand{\eps}{\varepsilon}
\newcommand{\dd}{\,\mathrm{d}}
\begin{document}
\section*{UnsolvedMath ID: 30005572}
\textbf{Source problem code:} catalogue entry
\subsection{Definitions and mathematical statement}
Write \(\widetilde X_1(N)\), \(\mathbb Z'\), and \(\mathcal O(N)\) for the space, coefficient ring, and arithmetic ring intended in the supplied report extract, and put
\[
A_N=H_1(\widetilde X_1(N),\mathbb Z'),\qquad
T_N=K_2(\mathcal O(N))\otimes\mathbb Z'.
\]
Here \(H_1\) denotes first homology. The extract does not define these objects, the admissible levels \(N\), the tensor-product base ring, or the convention for \(K_2\).

Let \(H_N^{\mathrm{rel}}\) denote the intended homology group relative to the non-infinity cusps, and let
\[
j_N:A_N\longrightarrow H_N^{\mathrm{rel}}
\]
be the comparison map implicit in the report's restriction language. The relative group, its ambient space, and its cusp subset require specification from the source context.

\textbf{Question, conditional on those specifications.} Does there exist a homomorphism
\[
\widetilde\Pi:H_N^{\mathrm{rel}}\longrightarrow T_N
\]
such that
\[
\Pi=\widetilde\Pi\circ j_N:A_N\longrightarrow T_N
\]
is Eisenstein and, for every Cremona symbol \([u:v]\) in the intended relative presentation,
\[
\widetilde\Pi([u:v])
 =\{\theta_E(cP),\theta_E(dP)\}?
\]
The braces denote the Steinberg symbol intended by the source. The extract does not define the Cremona symbols, their admissibility conditions or relations, the association of \(c,d\) to \(u,v\), the objects \(E,P,\theta_E\), or the hypotheses ensuring that the displayed symbol belongs to \(T_N\). It also supplies no Hecke operators or identities defining the Eisenstein condition. Consequently, this is a conditional transcription of the target, not a fully specified quantified conjecture; in particular, no range for \(N\) or additional localization is imposed here.

The extension direction follows the original extract: the absolute-homology map \(\Pi\) is the restriction, through \(j_N\), of a relative-homology map \(\widetilde\Pi\). The report also describes work in progress giving a basically affirmative answer, without stating its precise qualifications in the supplied passage.
\subsection{Short English statement}
The question asks whether a relative Bianchi homology map can send Cremona symbols to the displayed Steinberg symbols while inducing an Eisenstein map on absolute homology. Essential notation, level hypotheses, and the Eisenstein condition are missing from the extract, so the statement remains conditional and requires human mathematical review.
\subsection{Sources}
\begin{itemize}
\item Supplied local UnsolvedMath registry record, permanent ID \(30005572\), title \emph{Eisenstein Maps from Bianchi Homology to K2}, source problem code \texttt{catalogue entry}, in the Oberwolfach Reports Open Problems collection. The record gives proposed year 2023 and no named proposer.
\item \emph{Algebraische Zahlentheorie} (2023), Oberwolfach Reports, DOI 10.4171/owr/2023/28, pp.\ 1543--1606. Supplied report reference: \url{https://doi.org/10.4171/owr/2023/28}. The supplied background describes these reports as records of mathematical workshops and their open-problem discussions.
\item Sharifi, \emph{Eisenstein cocycles for imaginary quadratic fields}. Supplied literature reference: \url{https://doi.org/10.4171/owr/2023/28}.
\item Lecouturier--Sharifi--Shih--Wang, \emph{Eisenstein cocycles for imaginary quadratic fields}. Supplied literature reference: \url{https://arxiv.org/abs/2504.19125}.
\item Literature assessment dated exactly \texttt{2026-08-22}: the supplied curation labels the problem partially solved and states that the 2025 primary paper constructs the Bianchi-space to \(K_2\) map and establishes its Eisenstein property away from the level after controlled localization, while the exact integral all-level formulation in the extracted question is not fully established. This is supplied dated curation, not an independently verified assessment; the background expressly requests independent verification of time-sensitive status and cited claims.
\item The registry labels the statement \texttt{corrected\_verified} and reports that it was checked against the official Oberwolfach report. That verification assertion is supplied metadata, not a source check performed for this draft.
\end{itemize}
\end{document}
